\chapter{Depression in Children}

\section*{Definition and Overview}
Depression in children is a persistent low mood, loss of interest, and reduced enjoyment that affects a child's functioning at home, school, and with friends. It differs from the normal moodiness and sadness of childhood, which pass quickly, in its severity, duration, and impact. Depression in children is real, common, and treatable, but it is often missed because children may not describe their feelings in words and instead show irritability, behavioural change, or physical symptoms. This chapter explains the signs of depression in children, how it is diagnosed, and how it is treated.

\section*{Epidemiology}
Depression is one of the most common mental health problems in young people. Around one in twenty children and adolescents experiences depression, with rates rising sharply through the teenage years. Before puberty, depression affects boys and girls similarly; after puberty, it becomes more common in girls. Around half of adult mental health disorders begin before age 14, making early recognition and treatment vital. Rates of depression in Australian young people have increased in recent years, with social pressures, bullying, and family factors contributing.

\section*{Aetiology and Pathophysiology}
\begin{itemize}
    \item \textbf{Genetic factors:} a family history of depression increases risk
    \item \textbf{Brain and biology:} neurotransmitter and stress hormone changes affect mood regulation
    \item \textbf{Family factors:} parental depression, conflict, and family stress contribute
    \item \textbf{Adverse experiences:} bullying, abuse, neglect, and loss are major risk factors
    \item \textbf{Social factors:} peer problems, school pressure, and social media use
    \item \textbf{Personality and coping:} low self-esteem and poor coping skills increase vulnerability
    \item \textbf{Physical health:} chronic illness and some medications can contribute
\end{itemize}

\section*{Clinical Features}
\begin{itemize}
    \item Persistent sadness, tearfulness, or irritability and anger
    \item Loss of interest in activities, hobbies, and friends
    \item Withdrawal from family and social activities
    \item Changes in sleep: difficulty sleeping, nightmares, or sleeping too much
    \item Changes in appetite and weight
    \item Poor concentration, falling school performance, and school refusal
    \item Physical complaints such as headaches and stomach aches
    \item Low self-esteem, excessive guilt, and self-blame
    \item Thoughts of death or self-harm, which require immediate attention
    \item Symptoms persist for most of the day, most days, for more than two weeks
\end{itemize}

\section*{Differential Diagnosis}
\begin{itemize}
    \item Normal moodiness and emotional reactions to life events, which are shorter and less impairing
    \item Anxiety disorders, which often coexist with depression
    \item Attention deficit hyperactivity disorder, which can look like inattention and irritability
    \item Behavioural disorders such as oppositional defiant disorder
    \item Grief after a loss, which is time-limited
    \item Physical illness, including thyroid disorders and anaemia
    \item Autism spectrum conditions with mood difficulties
\end{itemize}

\section*{When to Seek Medical Care}
Parents should see a doctor if a child's low mood or behavioural change persists, affects school or friendships, or is accompanied by physical symptoms. Any mention of self-harm, hopelessness, or thoughts of death requires urgent assessment. School counsellors, GPs, and child mental health services all provide pathways to help.

\textbf{Call 000 (or local emergency services) for:}
\begin{itemize}
    \item Any immediate risk of self-harm or suicide
    \item A child who has harmed themselves or made plans to do so
    \item Severe agitation, confusion, or inability to keep the child safe
\end{itemize}
\textbf{Support lines:} Kids Helpline (1800 551 800) and Lifeline (13 11 14) are available 24 hours.

\section*{Diagnosis and Investigations}
\begin{itemize}
    \item \textbf{Clinical assessment:} interviews with the child and parents covering mood, behaviour, and functioning
    \item \textbf{Screening tools:} questionnaires such as the Patient Health Questionnaire for Adolescents
    \item \textbf{Information from school:} reports of behaviour and performance from teachers
    \item \textbf{Risk assessment:} careful enquiry about self-harm and suicidal thoughts
    \item \textbf{Physical assessment:} to exclude medical causes
    \item \textbf{Specialist referral:} to a child and adolescent psychiatrist or psychologist for complex cases
\end{itemize}

\section*{Management}
\begin{itemize}
    \item \textbf{Psychological therapy:} cognitive-behavioural therapy and interpersonal therapy are the first-line treatments for children and adolescents
    \item \textbf{Family involvement:} parenting support and family therapy address the family context
    \item \textbf{School support:} working with the school to adjust expectations and provide support
    \item \textbf{Medication:} antidepressants, particularly fluoxetine, may be used for moderate to severe depression in adolescents under specialist supervision
    \item \textbf{Lifestyle:} regular sleep, physical activity, and healthy eating
    \item \textbf{Addressing contributing factors:} managing bullying, family conflict, and social pressures
    \item \textbf{Monitoring:} regular review to assess response, safety, and side effects
    \item \textbf{Crisis care:} hospital care when there is significant risk
\end{itemize}

\section*{Complications}
\begin{itemize}
    \item Self-harm and suicide, the most serious risks
    \item School failure and dropout
    \item Social isolation and damaged friendships
    \item Substance use as a coping strategy
    \item Recurrent depression in adolescence and adulthood
    \item Physical health problems linked to inactivity and poor sleep
    \item Family conflict and strain
\end{itemize}

\section*{Prognosis and Outcomes}
With early recognition and treatment, most children and adolescents with depression recover. Psychological therapy is effective in the majority, and outcomes are better when treatment is started early and the family is involved. Untreated depression can last many months and recur, so ongoing support and monitoring are important. With good treatment and a supportive environment, most young people return to normal functioning and development.

\section*{Prevention and Self-Care}
\begin{itemize}
    \item Maintain open communication with children about their feelings
    \item Encourage healthy sleep, activity, and screen-time habits
    \item Build supportive family routines and time together
    \item Address bullying promptly through the school
    \item Seek help early rather than waiting to see if it passes
    \item Look after your own mental health as a parent
    \item For teenagers: encourage connection with friends and trusted adults, and limit social media pressure
\end{itemize}

\section*{Sources and Further Reading}
The following reputable sources provide further information for healthcare students and patients seeking reliable, current advice about depression in children.

\begin{itemize}
    \item Kids Helpline. \textit{Support and counselling for children and teenagers.} https://kidshelpline.com.au/
    \item Beyond Blue. \textit{Depression in children and teenagers.} https://www.beyondblue.org.au/
    \item headspace. \textit{Mental health support for young people.} https://headspace.org.au/
    \item Raising Children Network. \textit{Depression in children: signs and support.} https://raisingchildren.net.au/
    \item NHS. \textit{Depression in children and young people.} https://www.nhs.uk/mental-health/children-and-young-adults/
    \item Mayo Clinic. \textit{Depression in children: symptoms and treatment.} https://www.mayoclinic.org/
\end{itemize}
